\chapter{从开关到可调时钟：Science（2008）中的正负反馈}
前面的 R8 例子已经说明，正反馈不只用于记忆，还能与负反馈一起产生振荡。但“有振荡”与“适合做时钟”仍是两个不同的问题：改变周期时，峰值会不会一起消失？参数有一些变化时，振荡是否还能维持？

Tsai、Choi、Ma、Pomerening、汤超与 Ferrell 的 \emph{Robust, Tunable Biological Oscillations from Interlinked Positive and Negative Feedback Loops}（Science, 2008，参考文献~8）研究了这一问题。他们在所比较的细胞周期及简化振荡模型中发现，正负反馈交织可支持较宽的频率调节范围，同时维持近似不变的振幅；随机参数比较也显示了较大的振荡参数区域。这是模型之间的性能比较，不是“任意网络加一条正边都会更好”的定理。

\section{评价振荡：周期、振幅与参数范围要分别看}
设相邻两个峰的时间间隔为 $T$，频率为 $f=1/T$，一周期内的最大、最小浓度为 $x_{\max},x_{\min}$，振幅在本文记为
\[
A_{\rm osc}=x_{\max}-x_{\min}.
\]
比较振荡器时，至少要分开三个问题：它能否持续振荡，频率能调多少，以及调频时振幅改变多少。一个只能在很窄的参数区间里产生微弱摆动的系统，与一个可以调节周期、又保持明确高低状态的系统，虽然都能画出周期轨迹，功能上并不相同。

也要说明调的是哪个参数。如果把\emph{所有}速率统一乘以同一系数，等于把整个时间坐标一起缩放，当然可以改变周期而保留轨迹形状。更有意义的问题是，只改变一个合适的产生率或时间尺度，其他过程不作同步调整时，是否仍能稳定调频。稳健性则是另一个量：它考察某个采样范围内，多少组参数产生符合标准的振荡；数值取决于模型、采样规则和振荡判据。

\section{迟滞开关与慢负反馈，怎样把“高度”和“等待时间”分开？}
论文以细胞周期模型说明这种组合：CDK1 的正反馈包括激活 Cdc25、抑制 Wee1；经 APC 作用的负反馈则帮助结束高活动阶段。正反馈使切换更明确，负反馈使系统可以重新进入下一轮，而不是一直停在高状态（参考文献~8）。

为了把这个解释接到现有方程上，继续看本文自己的两变量模型，而不把它当作论文细胞周期模型的复现：
\[
\dot x=f(x,y),\qquad\dot y=\varepsilon g(x,y),
\]
其中 $f$ 与 $g$ 就是上一章的 R8 产生--降解规则。$\varepsilon$ 较小时，$x$ 快，抑制变量 $y$ 慢。暂时固定 $y$，只让 $x$ 恢复，就能画出\emph{快变量稳态曲线} $f(x,y)=0$。如果某个 $y$ 区间内存在明显分开的低表达和高表达两个稳定分支，中间夹着不稳定分支，$x$ 会在慢变量推动下沿一个稳定分支走，到折叠点才跳到另一支。

这就产生\emph{迟滞}：从低态转向高态与从高态转向低态，发生在不同的 $y$ 阈值。$y$ 增大时，高表达态一直保持到其分支终止，再快速降到低态；$y$ 减小时，低态又保持到另一个阈值，再快速升高。慢变化与快跳跃轮流发生，称为\emph{弛豫振荡}。这里“弛豫”是这一类振荡的标准名称，不表示整个过程都以单个指数恢复。

设高分支上 $y$ 从 $y_-$ 增到 $y_+$，低分支上再从 $y_+$ 降到 $y_-$，快分支上的稳态值分别记为 $x_H(y),x_L(y)$。若快慢分离足够明显、快跳跃所占时间很短，那么由 $\dd t=\dd y/[\varepsilon g(x,y)]$ 得到周期的近似：
\[
T\simeq\frac1\varepsilon\left[
\int_{y_-}^{y_+}\frac{\dd y}{g(x_H(y),y)}+
\int_{y_+}^{y_-}\frac{\dd y}{g(x_L(y),y)}\right].
\]
高分支上 $g>0$、低分支上 $g<0$，所以两项时间都为正。这说明慢变量的速度主要控制等待多久，而分支位置与跳跃阈值主要控制浓度跨越的范围。只改变 $\varepsilon$ 时，快稳态曲线本身不变，因此在这一近似适用的区间内，有可能明显改变周期而只小幅改变振幅。

\begin{figure}[H]\centering
\includegraphics[width=\linewidth]{figures/positive_feedback_tuning.pdf}
\caption{本文 R8 模型的快慢机制演示。左：快变量稳态曲线的稳定分支、不稳定分支与数值轨迹；中：改变同一个慢时间尺度，周期不同而浓度跨度接近；右：频率与振幅相对参考参数的变化。固定 $a_0=0.4,b_0=4,d_0=2$，仅改变 $\varepsilon$。本图不是 Science 论文原模型的数据。}\label{fig:positive_tuning}
\end{figure}

这个近似有明确限度。快慢分离减弱、参数靠近分支消失的位置，或反馈改变了快稳态曲线本身时，振幅与周期仍可能一起变化。新增脚本在所示参数范围内检查了后期重复峰、正浓度与频率变化；这些数值检查说明例子怎样工作，并不替代极限环存在性的证明。

\section{与 Cell 适应论文连接：正反馈的功能取决于回路角色}
这一讨论把正反馈的三种角色串起来：单独形成足够强的放大时，它可以提供多个稳定状态；与较慢负反馈配合时，它可以让系统在状态之间反复切换；在适应模型的特殊速率条件下，它还可能参与累计误差的调节。

Science（2008）关注的是时钟的可调性与振荡稳健性，后面 Cell（2009）关注的是响应之后能否回到基线。相似的正、负边可能被这两种功能以不同方式使用，评价标准也不同。因此，增加第三个节点时，我们不仅问是否更容易振荡，也问它是否能提供时间筛选、比例抵消或误差缓冲等新的动态角色。
